\chapter{Contexte général du projet}
\thispagestyle{MyStyle}
\begin{onehalfspace}

\section{Introduction}

Dans un contexte où la transformation numérique redéfinit les pratiques de l'administration publique, la gestion des concours de recrutement constitue un enjeu majeur. Les organismes publics sont confrontés à un volume important de candidatures et à la nécessité d'organiser, dans des délais contraints, l'affectation des candidats aux centres, établissements et salles d'examen, puis l'envoi des convocations. Un traitement encore largement manuel expose à des erreurs, des retards et une charge de travail élevée pour les gestionnaires. Le présent projet s'inscrit dans cette dynamique en proposant une plateforme de dématérialisation capable d'automatiser, en particulier, la répartition des candidats et de fiabiliser l'ensemble du processus, de l'import des candidatures jusqu'à l'envoi des convocations. Ce chapitre présente le cadre général du projet, l'organisme d'accueil, ainsi que les problématiques identifiées, la méthodologie retenue et la planification des travaux.

\section{Organisme d'accueil}

\subsection{Présentation générale}

\begin{figure}[H]
  \centering
  
\includegraphics[width=0.62\textwidth]{mef-logo.png}
  \caption{Logo du Ministère de l'Économie et des Finances}
  \label{fig:logo-mef}
\end{figure}

Le Ministère de l'Économie et des Finances (MEF) est le département ministériel du gouvernement du Royaume du Maroc chargé d'élaborer et de mettre en œuvre la politique économique et financière de l'État. Son siège est situé au boulevard Mohammed~V, quartier administratif, Rabat-Chellah.

Acteur central de la gouvernance publique, le MEF intervient dans la préparation et l'exécution de la loi de finances, la collecte des recettes, le paiement des dépenses, le contrôle des finances publiques, ainsi que dans la tutelle financière des établissements et entreprises publics. Il dispose d'une administration centrale et d'un réseau de services déconcentrés, et emploie un effectif important : plus de 17\,000 agents dans les directions à réseau et plus de 2\,000 dans les directions centrales, selon les données publiées par le Ministère.

Le MEF s'inscrit également dans une dynamique de modernisation de l'administration, notamment à travers la digitalisation de ses métiers et la gestion de son capital humain. C'est dans ce cadre que s'inscrit le présent stage, consacré à la dématérialisation de la gestion des candidatures aux concours de recrutement.

\subsection{Missions}

Conformément à ses attributions officielles, le Ministère est en charge des questions financières et monétaires, y compris les politiques de crédit et de finances extérieures. Il assure également la tutelle financière des entreprises et établissements publics, et contribue, le cas échéant, à leur transfert au secteur privé dans le cadre du processus de privatisation.

Ses principales missions peuvent être regroupées comme suit :

\begin{itemize}
  \item \textbf{En matière budgétaire et fiscale :} élaboration de la loi de finances et suivi de son exécution ; définition et mise en œuvre des politiques fiscale et douanière ; recouvrement des recettes publiques et paiement des dépenses publiques.
  \item \textbf{En matière financière :} définition des conditions des équilibres financiers interne et externe ; tutelle réglementaire et technique des établissements et intermédiaires financiers ; représentation du Maroc au sein des instances financières régionales et internationales.
  \item \textbf{En matière de contrôle :} contrôle des recettes et des dépenses publiques ; contrôle financier des entreprises et établissements publics ; contrôle des finances des collectivités territoriales ; réglementation et contrôle de l'activité des compagnies d'assurance et de réassurance.
  \item \textbf{En matière de restructuration et de cession d'actifs :} audits des entreprises et établissements publics ; établissement de la liste des entités à céder ou à privatiser ; préparation des dossiers de cession et des contrats-programmes.
\end{itemize}

Parallèlement : gestion des biens mobiliers et immobiliers de l'État et défense des administrations publiques en justice, notamment via l'Agence judiciaire du Royaume.

\subsection{Organisation}

Le Ministère comprend, outre le cabinet du ministre, une administration centrale et des services déconcentrés. Ses principales directions et structures sont les suivantes :

\begin{itemize}
  \item l'Inspection générale des finances (IGF) ;
  \item l'Administration des douanes et impôts indirects (ADII) ;
  \item la Trésorerie générale du Royaume (TGR) ;
  \item la Direction générale des impôts (DGI) ;
  \item la Direction du Budget ;
  \item la Direction du Trésor et des Finances extérieures ;
  \item la Direction des Entreprises publiques et de la Privatisation ;
  \item la Direction des Domaines de l'État ;
  \item la Direction des Affaires administratives et générales (DAAG) ;
  \item la Direction des Études et des Prévisions financières ;
  \item l'Agence judiciaire du Royaume ;
  \item la Direction de la Concurrence, des Prix et de la Compensation ;
  \item la Direction de la Gouvernance de la Sécurité des Systèmes d'Information (DGSSI).
\end{itemize}

Le Ministère exerce également une tutelle sur plusieurs organismes, parmi lesquels l'Office des changes, la Caisse marocaine des retraites (CMR) et Tamwilcom (Société nationale de garantie et de financement de l'entreprise).

\subsection{Direction d'accueil}

Le stage a été effectué au sein de la \textbf{Direction des Affaires administratives et générales (DAAG)}. Cette direction est notamment chargée de définir et de mettre en œuvre la politique de gestion des ressources humaines, de coordonner le recrutement, le déroulement de carrière et l'action sociale, ainsi que de promouvoir les systèmes d'information et les technologies de l'information au service des métiers du Ministère.

La DAAG constitue donc le cadre naturel du projet : l'organisation des concours de recrutement, l'affectation des candidats aux centres et salles d'examen, puis l'envoi des convocations relèvent directement de la gestion des ressources humaines et de la modernisation des processus administratifs. La plateforme développée vise à assister les gestionnaires dans ces opérations, en automatisant en particulier la répartition des candidats.

\section{Problématique et solutions proposées}

\subsection{Problématique}

L'organisation d'un concours au Ministère de l'Économie et des Finances ne s'arrête pas à la sélection des admis. Une fois la liste des candidats établie, les gestionnaires doivent encore affecter chaque candidat à un centre, un établissement et une salle, puis produire et transmettre les convocations. Ces opérations restent souvent réalisées de manière manuelle ou semi-manuelle (tableurs, échanges de fichiers, envois individuels), ce qui pose plusieurs difficultés.

D'abord, le volume de candidatures et la contrainte de délai rendent l'affectation aux salles longue et sujette aux erreurs : dépassement de capacité, double affectation, oubli d'un candidat ou affectation dans un centre éloigné de sa ville de résidence. Ensuite, l'absence d'une logique unique de répartition empêche de garantir une affectation équitable et cohérente, notamment selon la proximité géographique. Enfin, la génération et l'envoi des convocations restent fragmentés : risques de retard, d'erreur d'adresse ou d'information incomplète (salle, place, date et heure). Les responsables ne disposent pas non plus d'une vision consolidée (remplissage des salles, candidats non affectés, historique des envois).

La problématique peut donc être formulée ainsi : \textit{comment concevoir une plateforme permettant de gérer les candidatures aux concours, en automatisant la répartition des candidats dans les salles et en fiabilisant l'émission des convocations, tout en offrant aux gestionnaires un suivi opérationnel clair~?}

\subsection{Solutions proposées}

Pour répondre à cette problématique, le projet propose une application web dédiée à la gestion des concours, structurée autour des besoins des gestionnaires et des administrateurs.

La solution s'articule autour des axes suivants :

\begin{itemize}
  \item \textbf{Centralisation des données métier :} gestion des concours, des centres, des établissements et des salles, ainsi que l'import des candidats admis (fichier Excel), avec possibilité de modification, de suppression ou d'affectation manuelle.
  \item \textbf{Répartition automatique :} un moteur d'affectation place les candidats dans les salles selon le concours, la capacité disponible et la proximité géographique (ville / région), tout en signalant les cas non affectés (capacité insuffisante, ville non identifiée, absence de centre).
  \item \textbf{Convocations dématérialisées :} génération des convocations au format PDF (identité du candidat, concours, date et heure, centre, établissement, salle et numéro de place), puis envoi groupé par courrier électronique, avec un historique des envois.
  \item \textbf{Pilotage opérationnel :} un tableau de bord présente les indicateurs essentiels (effectifs, remplissage des salles, dernière répartition, suivi des convocations).
  \item \textbf{Contrôle d'accès :} distinction entre l'administrateur (supervision et gestion des comptes gestionnaires) et le gestionnaire (pilotage opérationnel du processus).
  \item \textbf{Architecture évolutive :} découpage en microservices (authentification, candidats, concours, lieux, répartition, convocations), afin d'isoler les responsabilités et de faciliter la maintenance.
\end{itemize}

Cette approche vise à réduire le travail répétitif, à limiter les erreurs d'affectation et à accélérer la notification des candidats, tout en conservant un contrôle humain (répartition manuelle, aperçu des convocations avant envoi).

\section{Objectifs du projet}

\subsection{Objectifs généraux}

L'objectif général de ce projet est de concevoir et de réaliser une plateforme de dématérialisation de la gestion des candidatures aux concours du Ministère de l'Économie et des Finances, afin d'assister les gestionnaires dans l'organisation des épreuves, d'automatiser la répartition des candidats dans les salles et de fiabiliser l'envoi des convocations.

\subsection{Objectifs spécifiques}

Cet objectif se décline en plusieurs objectifs spécifiques :

\begin{itemize}
  \item \textbf{Gérer le référentiel des concours et des lieux d'examen :} permettre la création, la modification et la suppression des concours, des centres, des établissements et des salles, en associant à chaque salle un concours et une capacité.
  \item \textbf{Dématérialiser le traitement des candidatures :} importer la liste des candidats admis à partir d'un fichier Excel, puis autoriser la consultation, la mise à jour, la suppression et, si besoin, l'affectation manuelle d'un candidat à une salle.
  \item \textbf{Automatiser la répartition :} affecter les candidats aux salles selon le concours, la capacité disponible et la proximité géographique (ville / région), tout en signalant les cas non affectés.
  \item \textbf{Générer et envoyer les convocations :} produire les convocations au format PDF (identité, concours, date et heure, lieu et numéro de place) et les transmettre par courrier électronique, avec un suivi des envois.
  \item \textbf{Offrir un pilotage opérationnel :} mettre à disposition un tableau de bord présentant les indicateurs essentiels (effectifs, remplissage des salles, dernière répartition, historique des convocations).
  \item \textbf{Sécuriser l'accès :} distinguer les rôles administrateur (supervision et gestion des comptes) et gestionnaire (pilotage opérationnel), via une authentification adaptée.
  \item \textbf{Adopter une architecture maintenable :} structurer la solution en microservices afin d'isoler les responsabilités (authentification, candidats, concours, lieux, répartition, convocations) et de faciliter l'évolution du système.
\end{itemize}

\section{Méthodologie de travail}

\subsection{Méthode en cascade utilisée}

\begin{figure}[H]
  \centering
  \includegraphics[width=0.92\textwidth]{modele-cascade.png}
  \caption{Cycle de développement selon la méthode en cascade}
  \label{fig:method-cascade}
\end{figure}

La conduite du projet a suivi une approche séquentielle, de type cycle en cascade (waterfall). Les besoins ayant été formalisés dès le départ dans le cahier des charges, il a été jugé préférable de stabiliser la conception avant le développement, plutôt que d'avancer par sprints courts. Le travail a été réalisé seul, avec l'accompagnement de l'encadrant : chaque tâche était présentée pour validation ; en cas d'accord, le projet passait à l'étape suivante ; en cas de réserve, les corrections étaient apportées avant de continuer.

Cette organisation combine deux niveaux. Au niveau global, les phases s'enchaînent dans un ordre fixe. Au niveau de chaque livrable, la revue avec l'encadrant joue le rôle de contrôle qualité, afin de limiter les écarts par rapport aux besoins exprimés.

Le déroulement du projet s'articule autour de quatre phases principales :

\begin{itemize}
  \item \textbf{Cahier des charges :} recueil et formalisation des besoins fonctionnels et non fonctionnels (acteurs, parcours gestionnaire, import des candidats, répartition, convocations, rôles et contraintes techniques).
  \item \textbf{Conception :} élaboration des diagrammes UML (cas d'utilisation, diagrammes de séquence et diagrammes de classes) afin de définir le comportement du système, les interactions et le modèle métier avant l'écriture du code.
  \item \textbf{Réalisation :} implémentation de la solution (interface, services métier, base de données), en respectant la conception validée, puis ajustements suite aux retours de l'encadrant.
  \item \textbf{Tests :} vérifications, corrections demandées par l'encadrant et contrôle du bon fonctionnement de l'ensemble.
\end{itemize}

\subsection{Répartition des tâches}

Les responsabilités ont été réparties de manière claire entre la conception et la réalisation technique, assurées par moi-même, et le suivi fonctionnel ainsi que la validation métier, assurés par l'encadrant. La partie technique a couvert l'analyse des besoins, la conception UML, le développement du backend (microservices et APIs), la réalisation des interfaces, la logique de répartition automatique, la génération et l'envoi des convocations, ainsi que les tests. De son côté, l'encadrant est intervenu dans la définition des besoins, la validation des livrables à l'issue de chaque tâche et le pilotage global du projet.

Cette organisation simplifiée a permis une communication directe, une bonne réactivité dans la prise de décision et une progression continue du projet.

\section{Planification du projet}

Le stage s'est déroulé du 6 mars au 6 juin 2026, soit une durée de trois mois. La planification a respecté l'enchaînement séquentiel retenu avec l'encadrant : d'abord la formalisation du besoin, puis la conception, ensuite le développement, et enfin les tests.

Les deux premières semaines (6--19 mars) ont été consacrées à la discussion du cahier des charges, à la compréhension du projet et à la définition de la manière de le réaliser. La conception a ensuite duré un mois (20 mars -- 19 avril) : élaboration des diagrammes de cas d'utilisation, de séquence et de classes, validés avant le passage au code.

La réalisation s'est étendue sur cinq semaines (20 avril -- 24 mai) : développement de la plateforme (authentification, concours et lieux d'examen, import des candidats, répartition automatique, convocations et tableau de bord). La période restante (25 mai -- 6 juin) a été dédiée aux tests, aux corrections demandées par l'encadrant et à la vérification du bon fonctionnement de l'ensemble.

Le diagramme de Gantt ci-dessous synthétise ces quatre phases.

\begin{figure}[H]
  \centering
  \includegraphics[width=\textwidth]{gantt.png}
  \caption{Diagramme de Gantt --- Planification du projet}
  \label{fig:gantt}
\end{figure}

Cette planification a permis de structurer efficacement le projet, de coordonner les activités et de suivre l'avancement de manière progressive et maîtrisée.

\section{Conclusion}

Ce chapitre a présenté le cadre général du projet, ainsi que l'environnement professionnel dans lequel il s'inscrit, à savoir le Ministère de l'Économie et des Finances, à Rabat. Il a permis de mettre en évidence les limites d'une gestion encore largement manuelle des candidatures --- affectation aux salles, génération et envoi des convocations --- et de justifier la nécessité d'une plateforme de dématérialisation capable d'automatiser la répartition des candidats tout en conservant un contrôle humain.

Il a également clarifié les objectifs du projet, les solutions proposées, la démarche méthodologique adoptée (cycle en cascade, avec validation de chaque tâche par l'encadrant) ainsi que la planification globale du travail, visualisée par le diagramme de Gantt. Cette synthèse constitue une base essentielle pour la suite du rapport et introduit naturellement le chapitre suivant, consacré au cahier des charges ainsi qu'aux besoins fonctionnels et non fonctionnels.

\end{onehalfspace}
